\documentclass[11pt,a4paper]{article}
\usepackage[utf8]{inputenc}
\usepackage[british]{babel}
\usepackage{geometry}
\geometry{margin=1in}
\usepackage{graphicx}
\usepackage{booktabs}
\usepackage{longtable}
\usepackage{float}
\usepackage{amsmath}
\usepackage{amssymb}
\usepackage{hyperref}
\usepackage{caption}
\usepackage{array}
\usepackage{ragged2e}

\hypersetup{colorlinks=true, linkcolor=black, citecolor=black, urlcolor=blue}

\title{Risk Management \\ GWP3\\[0.3em]}
\author{Ge Zhang \\ Alberto Hernandez-Espinosa}
\date{\today}

\begin{document}

\maketitle



\vspace{0.6em}

\section*{On AI usage}

This research used artificial intelligence tools in three limited ways: first, to help with
literature search and to check academic sources and datasets; second, when refining
algorithmic implementations and their validation; third, for grammar,
spelling and style corrections. All intellectual contributions (the
theoretical framing, the empirical analysis and the final reading of the results) remain entirely our own.

% \tableofcontents
%\newpage

%==========================================================================
\section{Introduction}

This report closes the mini-capstone begun in GWP1 and GWP2. The study underneath it is still Danish A.
Alvi's \textit{Application of Probabilistic Graphical Models in
Forecasting Crude Oil Price} (University College London, 2018). That dissertation claims that the price
of crude oil can be forecast in three steps: first, by reducing a panel of macroeconomic and physical-market
time series into latent regimes via hidden Markov models; second, by learning a belief network over those
discretised variables with a hill-climbing search; and third, by performing probabilistic inference on
the network to recover the regime of the oil price in the coming month.

Forecasting the price of crude oil is a well-studied problem, and the literature already tells us
that its moves come from different types of demand and supply shock (Kilian; Baumeister and Kilian).

Our own findings and conclusions are gathered at the end of the report.

%==========================================================================
\section{Training, validation and testing sets}

Our dataset has 199 month-end observations, running from January 2010 to July 2026.
It is put together as follows.
\begin{itemize}
    \item \textbf{Variables}: Seven time series in total (see Figure~\ref{fig:panel_gauss}).
    Three come from Yahoo Finance: front-month WTI and Brent crude futures, plus the US Dollar Index.
    The remaining four come from the Federal Reserve Economic Data service (see United States, Federal
    Reserve Bank of St. Louis): the Consumer Price Index, the effective federal funds rate, the industrial
    production index, and the WTI spot price.
    \item \textbf{Temporal alignment}: Every series was snapped to the last business day of the month
    so that all seven variables line up in time.
    \item \textbf{Outlier handling}: Very large values were capped using the interquartile-range rule,
    with a multiplier of three, to stop outliers from pulling the rest of the analysis around.
    \item \textbf{Gap filling}: Where the aligned series had missing observations, they were filled in
    by linear interpolation along the time axis.
\end{itemize}

The dataset is then split chronologically into three windows. Table~\ref{tab:roles} sets this division out,
and the sub-sections below add detail on each window in turn.

\begin{table}[H]
\centering
\caption{The division of the dataset between the three windows.}
\label{tab:roles}
\renewcommand{\arraystretch}{1.2}
\begin{tabular}{@{}lp{5.1cm}p{4.6cm}r@{}}
\toprule
\textbf{Set} & \textbf{What is estimated on it} & \textbf{How often it may be consulted} &
\textbf{Months} \\
\midrule
Training   & Hidden Markov parameters, network structure, conditional probability tables &
Freely & 139 \\
Validation & Nothing & Freely, but only to choose between specifications already fitted & 30 \\
Testing    & Nothing & Once, after every choice has been frozen & 30 \\
\bottomrule
\end{tabular}
\end{table}

As Table~\ref{tab:roles} makes plain, the split is strictly chronological:
training comes before validation, and validation before testing. Nothing later is ever used to fit
anything earlier.

\subsection{The training set}

The training set is the only window where parameters are actually estimated. Three estimation
problems are solved on it, and they are summarised in Table~\ref{tab:training_tasks}.

\begin{table}[H]
\centering
\caption{Estimation tasks performed on the training set}
\label{tab:training_tasks}
\renewcommand{\arraystretch}{1.3}
\begin{tabular}{@{}p{3.5cm}p{\dimexpr\textwidth-3.5cm-2\tabcolsep\relax}@{}}
\toprule
\textbf{Estimation Task} & \textbf{Implementation Details} \\
\midrule
Three-state hidden Markov model parameters $(\Pi, A, B)$ (one per time series) & Fitted via the Baum-Welch algorithm (Rabiner); re-estimated from 10 random initialisations to avoid local optima, retaining the solution with the highest log-likelihood. \\
\midrule
Belief network structure & Searched via hill climbing over the discretised training dataset (Chickering). \\
\midrule
Conditional probability tables (for each node given its parents) & Estimated under a K2 prior (Cooper and Herskovits) to avoid zero probability assignments for unobserved (but not impossible) state configurations, which would occur with maximum-likelihood estimation given the small sample size. \\
\bottomrule
\end{tabular}
\end{table}

As Table~\ref{tab:roles} also shows, the training window runs from January 2010 to July 2021,
139 months, or 69.8 per cent of the sample. Critically, it contains the two structural episodes
that define the modern crude market: the collapse of late 2014, driven by the expansion of US
shale supply, and the demand destruction of 2020. A training window missing either of those two
events would leave the model without any experience of crisis dynamics.

\subsection{The validation set}

Nothing is estimated on the validation set. Its job is to choose between specifications that have
already been fitted on the training set. It is needed because a specification that is chosen only
by how well it reproduces the training data is partly chosen for how well it reproduces noise.
Four decisions were settled here:

\begin{enumerate}
\item the scoring function for the structure search. Three candidates were considered: the K2 score, the Bayesian Dirichlet
      equivalent uniform score (Heckerman, Geiger and Chickering), or the discrete Bayesian
      information criterion;
\item the maximum in-degree allowed in the search, set at 2, 3 or unrestricted; this controls how
      many parents a node may take on, and therefore how quickly the conditional probability
      tables become thinly populated;
\item whether the search is seeded with an expert graph that encodes standard energy-market
      economics, or started from an empty graph;
\item the decision threshold that turns the posterior distribution over regimes into an actual
      position.
\end{enumerate}

This mirrors the source paper, which experiments with ``different Bayesian Dirichlet scoring
functions such as the BDeu or K2 and the Bayesian Information Criterion'' (Alvi 41) and notes that
``the validation step is to adjust the model if the error is too high'' (Alvi 62). The window itself
runs from August 2021 to January 2024, thirty months covering the post-pandemic recovery and the
energy shock that followed the invasion of Ukraine. Eighteen configurations were compared on it
for each of the two specifications reported in Section~\ref{sec:results}.

\subsection{The testing set}

The testing set is used only to measure. It was consulted under three rules:

\begin{enumerate}
\item it is opened once, after discretisation, structure, parameters and decision threshold have
      all been frozen;
\item nothing is adjusted after the fact, because each fresh look would simply turn it into a
      second validation window;
\item the error it reports is the number quoted, with no after-the-fact qualification or selection
      among variants.
\end{enumerate}

The window runs from February 2024 to July 2026, again 30 months. This follows the spirit of the
target reference, which suggests at least thirty observations before a proportion can be read
with any confidence (Alvi 41). A 70:15:15 split was therefore adopted instead of the 80:10:10
split used in the source. With 199 months, an 80:10:10 split would leave only twenty months on
each side, and the exact binomial confidence interval around an accuracy of, say, 60 per cent
would then stretch from roughly 36 to 81 per cent, far too wide to support any firm conclusion.
The 70:15:15 split keeps the paper's logic while giving 30 months to each hold-out window; even
so, and as Section~\ref{sec:results} makes clear, the intervals remain wide.

%==========================================================================
\section{Allocation of the data and comparison of the hold-out sets}

Table~\ref{tab:split} gives the basic statistics for each window, and
Figure~\ref{fig:allocation} plots the division on top of the price series itself.

\begin{table}[H]
\centering
\caption{Allocation of the 199 monthly observations and descriptive statistics of the WTI spot
price in each window.}
\label{tab:split}
\begin{tabular}{@{}lrlrrrrr@{}}
\toprule
\textbf{Set} & \textbf{Months} & \textbf{Period} & \textbf{Share} & \textbf{Mean} &
\textbf{Mean} & \textbf{Vol.} & \textbf{Range} \\
 & & & (\%) & price & log ret. & log ret. & (USD) \\
\midrule
Training   & 139 & 2010-01 to 2021-07 & 69.8 & 69.34 & \phantom{-}0.0001 & 0.1231 & 19.2--113.4 \\
Validation &  30 & 2021-08 to 2024-01 & 15.1 & 83.61 & \phantom{-}0.0037 & 0.0950 & 66.1--114.4 \\
Testing    &  30 & 2024-02 to 2026-07 & 15.1 & 73.39 & -0.0000 & 0.1196 & 57.3--108.6 \\
\bottomrule
\end{tabular}
\end{table}

\begin{figure}[H]
\centering
\includegraphics[width=\textwidth]{figures/data_allocation.png}
\caption{Chronological allocation of the sample. The division respects the arrow of time: every
observation used for model selection postdates every observation used for estimation, and every
observation used for testing postdates both.}
\label{fig:allocation}
\end{figure}

\subsection{Are the two hold-out windows comparable?}

The question that matters here is whether the validation and testing windows describe the market
in roughly comparable terms. That was tested on the monthly log returns of the WTI spot price
using three standard tests:
\begin{itemize}
    \item a two-sample Kolmogorov-Smirnov test for equality of the full distributions,
    \item Levene's test for equality of variances,
    \item and Welch's test for equality of the sample means.
\end{itemize}
Table~\ref{tab:ks} gives the results.

\begin{table}[H]
\centering
\caption{Tests of distributional equality between the windows, monthly log returns of the WTI spot
price.}
\label{tab:ks}
\begin{tabular}{@{}lrrrr@{}}
\toprule
\textbf{Comparison} & \textbf{KS statistic} & \textbf{KS $p$} & \textbf{Levene $p$} &
\textbf{Welch $p$} \\
\midrule
Training vs validation & 0.144 & 0.647 & 0.921 & 0.860 \\
Training vs testing    & 0.123 & 0.812 & 0.972 & 0.996 \\
Validation vs testing  & 0.172 & 0.791 & 0.953 & 0.895 \\
\bottomrule
\end{tabular}
\end{table}

None of the tests rejects at the five per cent level. On a statistical reading, then, the three
windows look much the same in central tendency, spread and overall shape. That is reassuring,
because it suggests that choosing the model on the validation window is not the same as tuning it
on a market that behaves completely differently from the one used for final testing. Even so, two
qualifications are essential.

First, failing to reject is not the same as proving the windows are equal. With around thirty
observations these tests have low power, and only a very large difference would reliably be
picked up. Second, similarity in the distribution of returns is not the same thing as similarity
in the underlying economic environment. The validation window is shaped by a war-driven supply
shock; the testing window is a period of gradual normalisation. Those two situations can easily
produce return distributions that look alike on paper even if they are driven by very different
mechanisms, and it is precisely the mechanism that a belief network claims to capture. The
comparison therefore licenses the validation procedure, but it does not guarantee that a
structure learned in one part of the sample will transfer cleanly to another.

\subsection{How the regimes are distributed within each window}

The distributional tests are run on the raw returns, but what the model actually
forecasts is the regime. The two hold-out windows therefore also need to be compared
in terms of their regime composition. Table~\ref{tab:shares} reports that breakdown
for the discretisation kept in Section~\ref{sec:regimes}.

\begin{table}[H]
\centering
\caption{Composition of each window by decoded regime of the WTI spot price, log-return emissions
(per cent of months).}
\label{tab:shares}
\begin{tabular}{@{}lrrrr@{}}
\toprule
\textbf{Set} & \textbf{Months} & \textbf{Bear} & \textbf{Stagnant} & \textbf{Bull} \\
\midrule
Training   & 138 & 21.7 & 65.9 & 12.3 \\
Validation &  30 & 23.3 & 63.3 & 13.3 \\
Testing    &  30 & 13.3 & 73.3 & 13.3 \\
\bottomrule
\end{tabular}
\end{table}

The training and validation windows look much alike, but the testing window is noticeably quieter
than either. The stagnant state accounts for 73.3 per cent of its months, against 63.3 per cent
in the validation window, and the bear share falls from 23.3 per cent to 13.3 per cent. The
practical consequences of this shift are summarised below and taken up again in
Section~\ref{sec:results}.

Table~\ref{tab:window_implications} pulls together the main differences between validation and
testing that matter when the model is being judged:

\begin{table}[H]
\centering
\caption{Key distinctions between validation and testing windows that impact model performance
assessment. The baselines are measured on the window itself; the 72.4 per cent reported in
Table~\ref{tab:accuracy} is the same rule measured over the 29 months of the one-month-ahead
target.}
\label{tab:window_implications}
\renewcommand{\arraystretch}{1.3}
\begin{tabular}{@{}p{3.5cm}p{5.2cm}p{5.2cm}@{}}
\toprule
\textbf{Metric} & \textbf{Validation Window} & \textbf{Testing Window} \\
\midrule
Always-stagnant baseline accuracy & 63.3\% & 73.3\% \\
Number of directional (non-stagnant) months & 11 (36.7\% of window) & 8 (26.7\% of window) \\
Total economically informative months for forecasting & 11 & 8 \\
\midrule
Core evaluation implication & Model accuracy must outperform the lower 63.3\% baseline to demonstrate skill & The higher 73.3\% baseline means raw accuracy can appear inflated without additional context; small sample of directional months produces wide confidence intervals that cannot be resolved by model refinement alone \\
\bottomrule
\end{tabular}
\end{table}

\subsection{Comparison by purpose}

Beyond the numbers, the two sets differ in what a poor result on each one would actually mean:

\begin{itemize}
    \item A poor result on the validation set is informative, and it is actionable. It tells us to
    change the scoring function, the in-degree limit or the discretisation, and we can do that as
    often as we like.
    \item A poor result on the testing set, by contrast, is terminal for the specification. To
    start re-working things at that stage would simply corrupt the final estimate.
\end{itemize}

This asymmetry is the reason both windows are used, and it is why the 18 configurations reported
in Section~\ref{sec:replication} were compared on the validation window alone.

%==========================================================================
\section{Regime detection: from prices to discrete evidence}
\label{sec:regimes}

A discrete belief network needs inputs drawn from a fixed and finite set of states, but our panel
is made of continuous series. Following the source paper, each series was therefore converted into
a sequence of three regimes (bear, stagnant, bull) using a three-state hidden Markov model
estimated on the training window alone (Rabiner). The states were then ordered by the arithmetic
mean of the underlying monthly change, so that the labels carry an economic meaning rather than
being arbitrary state numbers.

Two emission schemes were fitted and compared.

\paragraph{Parity emissions (the source paper's scheme).} Each series is first differenced, and each
change is then reduced to its sign. This yields a two-symbol alphabet: one symbol for an increase
and zero for a decrease (Alvi 42). The resulting model is a discrete-emission hidden Markov model
with three hidden states.

\paragraph{Log-return emissions (our modification).} In the modified scheme, the emission is the
monthly log return itself. To stop extreme observations from distorting things, returns are capped
at three training-sample standard deviations and then modelled with a Gaussian density with
diagonal covariance. The transition matrix receives a sticky Dirichlet prior, implemented by
adding pseudo-counts of five to the diagonal (Fox et al.). In practical terms, this encodes the
prior belief that market regimes tend to last for several months, rather than swapping every
month.

In both cases decoding is done \textit{in real time}. Concretely, the regime assigned to month
$t$ is obtained by running the Viterbi recursion on emissions up to and including month $t$ only.
This is a deliberate change from the source paper, which decodes each hold-out window as a whole
(Alvi 60, 63). Decoding an entire window at once lets the state assigned to a given month be
influenced by what happens later in the sample. The resulting regime series then contains
information that would not have been available at the time, and any forecast built on it is
contaminated by look-ahead bias. Our first implementation reproduced that choice and produced a
nominal one-month-ahead accuracy of 100 per cent, which is a tell-tale sign of leakage rather
than genuine predictive skill.

\subsection{What the two schemes produce}

Table~\ref{tab:hmm} summarises the fitted models for the WTI spot price, and
Figure~\ref{fig:transmat} plots the two transition matrices side by side.

\begin{table}[H]
\centering
\caption{Hidden Markov models of the WTI spot price estimated on the training window.}
\label{tab:hmm}
\begin{tabular}{@{}lrr@{}}
\toprule
 & \textbf{Parity emissions} & \textbf{Log-return emissions} \\
\midrule
Log-likelihood & $-92.37$ & $138.53$ \\
Average diagonal of the transition matrix & 0.000 & 0.742 \\
Regime switches over 198 months & 189 & 52 \\
Mean monthly log return, bear state & $-$ & $-0.129$ \\
Mean monthly log return, stagnant state & $-$ & $+0.014$ \\
Mean monthly log return, bull state & $-$ & $+0.152$ \\
\bottomrule
\end{tabular}
\end{table}

\begin{figure}[H]
\centering
\includegraphics[width=\textwidth]{figures/transition_matrices.png}
\caption{Estimated transition matrices for the WTI spot price. Under parity emissions the diagonal
is empty: the chain is obliged to leave its state at every step. Under log-return emissions the
mass concentrates on the diagonal, which is the behaviour a market regime is supposed to exhibit.}
\label{fig:transmat}
\end{figure}

The message is straightforward. The parity specification does not recover usable market regimes.
Because it sees only whether the price moved up or down, it produces a sequence that switches
state in 189 of 198 months and behaves almost like a simple alternation. The log-return
specification, by contrast, recovers persistent, economically interpretable states. Average
persistence is 0.74, which implies an expected duration of about four months, and the full sample
contains only 52 regime switches. This is behaviour much closer to what the regime-switching
literature describes (Hamilton; Ang and Timmermann). The log-return model is therefore the one
that is carried forward into the rest of the analysis. Table~\ref{tab:regime_summary} summarises
the comparison, Table~\ref{tab:latent} spells out the latent meaning of each state,
Figure~\ref{fig:regimes} plots the decoded sequence against the price, and
Figure~\ref{fig:regimes_parity} shows the same series under the parity scheme.

\begin{table}[H]
\centering
\caption{Why the log-return specification is more interpretable than the parity specification}
\label{tab:regime_summary}
{\small\setlength{\tabcolsep}{4pt}\renewcommand{\arraystretch}{1.15}\begin{tabular}{>{\raggedright\arraybackslash}p{2.8cm} >{\raggedright\arraybackslash}p{5.0cm} >{\raggedright\arraybackslash}p{5.0cm}}
\toprule
Feature & Parity emissions & Log-return emissions \\
\midrule
What the model sees & Only whether the monthly change is positive or negative. & The size and sign of the monthly log return, after capping extremes. \\
Decoded behaviour & Regimes switch in 189 of 198 months, so states alternate almost every month. & Regimes switch 52 times, which is much more consistent with persistent market phases. \\
Economic meaning & Weak: the states mainly separate signs, not genuine market conditions. & Strong: the states map naturally into bear, stagnant, and bull phases. \\
Persistence & Zero probability of remaining in the same state. & Average persistence is 0.74, implying a regime duration of about four months. \\
Overall interpretation & The model does not recover usable market regimes. & The model produces regimes an analyst can interpret and use. \\
\bottomrule
\end{tabular}}
\end{table}




\begin{table}[H]
\centering
\caption{Latent meaning of the hidden states under log-return emissions, training window.}
\label{tab:latent}
\begin{tabular}{@{}lrrrr@{}}
\toprule
\textbf{State} & \textbf{Months} & \textbf{Mean log return} & \textbf{Volatility} &
\textbf{Share (\%)} \\
\midrule
Bear     & 30 & $-0.1289$ & 0.1463 & 21.7 \\
Stagnant & 91 & $+0.0143$ & 0.0532 & 65.9 \\
Bull     & 17 & $+0.1518$ & 0.1350 & 12.3 \\
\bottomrule
\end{tabular}
\end{table}

\begin{figure}[H]
\centering
\includegraphics[width=\textwidth]{figures/regimes_full_gaussian.png}
\caption{Decoded regimes of the WTI spot price under log-return emissions. The bear state captures
the second half of 2014, the first months of 2020 and the 2025 correction; the bull state captures
the 2021 recovery and the first quarter of 2022.}
\label{fig:regimes}
\end{figure}

\begin{figure}[H]
\centering
\includegraphics[width=\textwidth]{figures/regimes_full_parity.png}
\caption{Decoded regimes of the same series under the parity emissions of the dissertation. The
colour changes almost every month; no episode of the past sixteen years is identified as a
sustained bear or bull market.}
\label{fig:regimes_parity}
\end{figure}

\subsection{The seven series, not only the price}

The belief network consumes the discretised version of all seven variables, so
Figures~\ref{fig:panel_gauss} and~\ref{fig:panel_parity} plot every series coloured by the regime
that its own hidden Markov model assigns to it. These labels are relative to each series, so for a
trending macroeconomic index ``bear'' means the state of slowest growth, not a falling market.

\begin{figure}[H]
\centering
\includegraphics[width=\textwidth]{figures/regimes_panel_gaussian.png}
\caption{Decoded regimes of all seven modelled series under log-return emissions. The two crude
futures and the spot price show blocks lasting several months; the Consumer Price Index, the
industrial production index and the dollar index are almost monochromatic.}
\label{fig:panel_gauss}
\end{figure}

\begin{figure}[H]
\centering
\includegraphics[width=\textwidth]{figures/regimes_panel_parity.png}
\caption{The same seven series under the parity emissions of the dissertation. The dollar index
changes regime in 197 of 198 months and the spot price in 189; Brent, by contrast, collapses onto a
single state. Neither behaviour is a market regime.}
\label{fig:panel_parity}
\end{figure}

The panel view reveals something that the single-series plot in Figure~\ref{fig:regimes} leaves
hidden. Under log-return emissions, the three oil-price series behave in a way that broadly fits
the idea of a regime process: WTI futures switch regime 46 times, Brent 35 times, and the spot
price 52 times over 198 months. The macroeconomic variables behave very differently. The Consumer
Price Index never leaves its first state, industrial production spends 97 per cent of the sample
in a single state, and the dollar index spends 88 per cent in one. At monthly frequency these
series are close to monotonic, so once they are modelled with Gaussian emissions and a prior that
favours persistence, one state absorbs virtually every observation. After discretisation the
variable is therefore left close to constant.

That matters because a nearly constant variable brings almost no information to the belief
network. It is for this reason that the structure search in Section~\ref{sec:results} drops the
Consumer Price Index, the federal funds rate and the dollar index while retaining the price-based
nodes. Put another way, the macroeconomic block of the model is silenced before the network is
even estimated. The bottleneck is the discretisation scheme, not the belief network itself. Any
future work on this pipeline would probably do better to discretise trending macroeconomic series
against a moving growth benchmark rather than on the raw monthly change.

Under parity emissions the failure takes the opposite shape, but it is just as serious. The dollar
index alternates in 197 of 198 months, while Brent, whose differenced monthly changes are
nearly always of the same sign, collapses onto a single state and switches only twice in
sixteen years. One series changes state almost every month, the other barely changes state at
all. Neither pattern corresponds to a useful economic idea of a regime.

This is the main economic reading of the two figures. A regime that lasts a single month is not
really a regime, and a system that reports a fresh market state each month cannot inform a
position meant to be held for a quarter. The advantage of the improved specification is therefore
not just a question of statistical tidiness. It changes the type of state sequence that the
downstream network is being asked to learn and predict.

%==========================================================================
\section{Re-running the Bayesian network with hill climbing}
\label{sec:replication}

The original paper's protocol was reproduced as literally as possible. A node called
\texttt{forecast}, which simply duplicates the discretised WTI spot-price regime, is added
to the training frame. The network structure is then searched by hill climbing, the conditional
probability tables are estimated under a K2 prior, and the network is asked for the most probable
state of \texttt{forecast} given the evidence in the seven remaining variables. The sequence of
inferences that comes out is then rolled forward by one month so that it can be read as a
forecast. This is the same one-month shift that the original paper applies in its own code (Alvi
62--63). The computations use open-source implementations of hidden Markov models and belief
networks (Ankan and Panda).

Eighteen configurations were compared on the validation window, crossing three scoring functions,
three in-degree limits and the presence or absence of the expert seed. The configuration retained
was the one using the K2 score with unrestricted in-degree, seeded with the expert graph. It
produced a network with 24 edges, in which the Markov blanket of the forecast node (Koller and
Friedman) contains all seven remaining variables. Figure~\ref{fig:network_k2} shows the network
learned by the unseeded K2 search on the training window.

\begin{figure}[H]
\centering
\includegraphics[width=0.85\textwidth]{figures/network_k2_train.png}
\caption{Belief network over the discretised variables learned by hill climbing with the K2 score
on the training window. The forecast node, shaded, is attached to the futures and macroeconomic
nodes, reproducing the densely connected topology reported by the dissertation (Alvi 60).}
\label{fig:network_k2}
\end{figure}

\subsection{Were the results of the paper replicated?}

Table~\ref{tab:replication} places our figures beside those reported in the dissertation.

\begin{table}[H]
\centering
\caption{Replication of the validation and testing errors of Alvi (2018, 62--63).}
\label{tab:replication}
\begin{tabular}{@{}llrrr@{}}
\toprule
\textbf{Study} & \textbf{Set} & \textbf{Months} & \textbf{Error (\%)} & \textbf{Accuracy (\%)} \\
\midrule
Alvi (2018)      & Validation & 28 & 67.86 & 32.14 \\
Alvi (2018)      & Testing    & 28 & 42.86 & 57.14 \\
This replication & Validation & 30 & 86.67 & 13.33 \\
This replication & Testing    & 30 & 76.67 & 23.33 \\
\bottomrule
\end{tabular}
\end{table}

The short answer is: \textbf{partially}.

What does replicate is the procedure, and its qualitative shape. The full pipeline runs end to end
on an independent dataset. The hill-climbing search returns a densely connected graph of the same
general type as the one shown in the original paper. The forecast node picks up the futures and
macroeconomic variables as parents, and the validation error is very high, just as it is in the
paper, where roughly two-thirds of the validation months are misclassified.

What does not replicate is the magnitude of the testing error. The original paper reports a
42.9 per cent error, equivalent to 57.1 per cent accuracy, and presents that figure as
encouraging. Our replication produces 76.7 per cent error, or 23.3 per cent accuracy. That is
\textit{worse} than the 33.3 per cent accuracy an uninformed guess would deliver if the three
states were equally likely. When a model performs consistently below chance, the problem is not
noise. It is evidence that the model is systematically wrong, and that kind of systematic error
needs an explanation.

The explanation lies in the interaction between the parity discretisation of
Section~\ref{sec:regimes} and the one-month rolling convention. The network only ever produces a
\textit{nowcast}: because the \texttt{forecast} node is an exact copy of the current month's
target regime, the structure search simply learns that deterministic relationship and replicates
the current regime almost perfectly. Rolling this perfect nowcast forward by one month turns its
accuracy into a measure of regime persistence; the forecast is only right if the regime happens
not to change.

Under the parity emission framework, regimes are constructed with near-zero persistence by design
(the estimated chain leaves its state at almost every step). The protocol therefore never really
measures the network's predictive power; it measures only how persistent the discretised labels
are. The 57 per cent testing accuracy in the original paper is plausibly explained by its sample
generating more persistent regimes from the same procedure. On our implementation, at least,
there is no evidence that the network genuinely forecasts future states.

This gives one of the central methodological conclusions of the exercise: if the forecast is to
come from the network, it has to be produced by the network itself, and not pulled out of the
persistence of the labelling scheme.

%==========================================================================
\section{Results: the accuracy of the crude oil forecast}
\label{sec:results}

Three specifications were used to judge the network's predictive power, all over the same months
of the testing window.

\begin{description}
\item[Replication model.] The original protocol as described above: parity emissions, a duplicated
forecast node, and the inferred state rolled forward by one month.

\item[Improved model.] The target is redefined as the regime of the \textit{following} month, so
the network is asked to forecast rather than reconstruct. The evidence is the current regime of
all seven series, under log-return emissions. Model selection is again performed on the validation
window. The retained specification uses the discrete Bayesian information criterion, a maximum
in-degree of two, and no expert seed. It produces a sparse four-edge network
(Figure~\ref{fig:network_pred}):
\texttt{WTI\_Spot}\,$\rightarrow$\,\texttt{WTI\_CL}\,$\rightarrow$\,\texttt{Brent\_BZ},
\texttt{WTI\_CL}\,$\rightarrow$\,\texttt{forecast}, and
\texttt{Brent\_BZ}\,$\rightarrow$\,\texttt{Ind\_Prod}. In that network the Markov blanket of the
forecast node reduces to the front-month WTI futures contract.

\item[Expected-return rule.] The posterior distribution over next month's regime, produced by the
improved model, is turned into a trading position by maximising the expected monthly log return
implied by that posterior, using the regime means from Table~\ref{tab:latent}. A threshold of one
per cent, chosen on the validation window, keeps the rule flat when the expected move is small.
\end{description}

\begin{figure}[H]
\centering
\includegraphics[width=0.85\textwidth]{figures/network_predictive.png}
\caption{The improved predictive network. Selection on the validation window favours a sparse
graph: with 138 training months and three states per node, dense networks exhaust the data in
their conditional probability tables.}
\label{fig:network_pred}
\end{figure}

\subsection{Classification accuracy}

Table~\ref{tab:accuracy} gives the headline numbers, and Figure~\ref{fig:bars} sets them against
three benchmarks: an uninformed guess, always predicting the modal regime from the training
window, and simply repeating last month's regime.

\begin{table}[H]
\centering
\caption{One-month-ahead forecast of the WTI spot regime. Balanced accuracy is the mean of the
per-regime recalls and is not inflated by the dominance of the stagnant state.}
\label{tab:accuracy}
\begin{tabular}{@{}llrrrr@{}}
\toprule
\textbf{Model} & \textbf{Set} & \textbf{Months} & \textbf{Accuracy} & \textbf{Balanced} &
\textbf{Macro} \\
 & & & (\%) & \textbf{accuracy (\%)} & \textbf{F1 (\%)} \\
\midrule
Replication          & Validation & 30 & 13.33 & 17.32 & 13.03 \\
Replication          & Testing    & 30 & 23.33 & 22.35 & 19.26 \\
Improved             & Validation & 29 & 58.62 & 34.39 & 30.95 \\
Improved             & Testing    & 29 & 75.86 & 41.67 & 39.31 \\
Expected-return rule & Testing    & 29 & 68.97 & 45.24 & 45.46 \\
\midrule
Uninformed guess     & Testing    & 29 & 34.48 & -- & -- \\
Majority regime      & Testing    & 29 & 72.41 & -- & -- \\
Persistence          & Testing    & 29 & 79.31 & -- & -- \\
\bottomrule
\end{tabular}
\end{table}

\begin{figure}[H]
\centering
\includegraphics[width=\textwidth]{figures/accuracy_comparison.png}
\caption{Accuracy on the testing window. The dashed line is the accuracy of an uninformed guess
over three states. Neither belief-network specification beats the persistence benchmark.}
\label{fig:bars}
\end{figure}

Three points stand out, and the third is the one that matters.

First, \textbf{the replication model performs worse than chance}. Its exact 95 per cent binomial
confidence interval runs from 9.9 to 42.3 per cent, for the structural reason already explained
in Section~\ref{sec:replication}.

Second, \textbf{the improved model classifies 75.9 per cent of months correctly, with an exact
interval from 56.5 to 89.7 per cent}. That is above the majority benchmark of 72.4 per cent, but
still below the persistence benchmark of 79.3 per cent. Exact McNemar tests on the paired
month-by-month outcomes (McNemar) give $p = 1.00$ against both benchmarks. Over a sample this
short, the belief network is therefore statistically indistinguishable from the very simple rule
of repeating last month's regime.

Third, \textbf{the balanced accuracy of the improved model is only 41.7 per cent}.
Figure~\ref{fig:conf} shows why. Under the maximum a posteriori rule, the model answers
``stagnant'' in twenty-six of the twenty-nine months. It correctly picks up the 21 stagnant
months, but it captures only one of the four bear months and none of the four bull months.
Figure~\ref{fig:timeline} tells the same story month by month: the model rides the long stagnant
stretch of 2024 and early 2025, and misses the turning points of mid-2025 and 2026. A rule that
is right three times out of four mainly because it keeps saying ``nothing much will happen'' may
look respectable on a headline accuracy measure, but it is not much use in practice. The
difference between raw accuracy and actual decision usefulness is exactly the distinction a risk
manager needs to keep in view.

\begin{figure}[H]
\centering
\begin{minipage}{0.49\textwidth}
\includegraphics[width=\textwidth]{figures/confusion_modelA_test.png}
\end{minipage}
\begin{minipage}{0.49\textwidth}
\includegraphics[width=\textwidth]{figures/confusion_modelB_test.png}
\end{minipage}
\caption{Confusion matrices on the testing window: replication model (left) and improved model
(right). The improved model concentrates its answers on the stagnant regime; the replication model
distributes its errors across all three.}
\label{fig:conf}
\end{figure}

\begin{figure}[H]
\centering
\includegraphics[width=\textwidth]{figures/timeline_modelB_test.png}
\caption{Month-by-month forecasts of the improved model on the testing window. Correct months are
marked with a circle and incorrect months with a cross. The model tracks the long stagnant stretch
of 2024 and early 2025 and misses the turning points of mid-2025 and 2026.}
\label{fig:timeline}
\end{figure}

\subsection{From probabilities to decisions}

The failure of the maximum a posteriori rule is a decision-theoretic failure, not a
probabilistic one. The posterior itself, shown in Figure~\ref{fig:posterior}, does move in
sensible ways: the probability of the bear state rises ahead of the mid-2025 correction, and the
probability of the bull state rises before the 2026 rebound, even though neither ever becomes the
modal state. Keeping only the most probable regime simply discards that information. The reason
is that the loss function implicit in the maximum a posteriori rule treats all three types of
error as equally costly, whereas in a commodity book the cost of missing a twelve per cent decline
is plainly not the same as the cost of missing a one per cent drift.

The expected-return rule puts the economics back in. Writing $p_t(k)$ for the posterior
probability of regime $k$ in month $t$ and $\mu_k$ for the mean monthly log return of that regime
estimated on the training window, the rule goes long when $\sum_k p_t(k)\mu_k$ exceeds a
threshold $\tau$, short when it falls below $-\tau$, and flat otherwise. The threshold is fixed
at $\tau = 0.01$ on the validation window.

The result, shown in Table~\ref{tab:accuracy}, is that raw accuracy falls from 75.9 to
69.0 per cent, while balanced accuracy rises from 41.7 to 45.2 per cent and macro F1 from 39.3 to
45.5 per cent. The rule issues a directional call in nine of the twenty-nine months, giving 31
per cent coverage, and it is right in 66.7 per cent of those calls.
Figure~\ref{fig:strategy} traces the value of one barrel traded on those signals: starting from
79.2 US dollars, it ends the testing window at 157.0 dollars, compared with 70.6 dollars for a
barrel held throughout.

\begin{figure}[H]
\centering
\includegraphics[width=\textwidth]{figures/posterior_modelC.png}
\caption{Posterior probability of each regime for the coming month over the testing window, with
the realised regime marked above. The bear probability rises before the mid-2025 decline without
ever becoming modal, which is what the maximum a posteriori rule discards.}
\label{fig:posterior}
\end{figure}

\begin{figure}[H]
\centering
\includegraphics[width=\textwidth]{figures/strategy_modelC.png}
\caption{Value of one barrel traded on the expected-return rule against a barrel held throughout
the testing window. Transaction costs, financing and slippage are excluded and the path rests on
nine decisions; the exhibit illustrates the decision rule and is not a performance claim.}
\label{fig:strategy}
\end{figure}

That backtest needs to be read with discipline. It excludes transaction costs, financing and
slippage; it assumes execution at the month-end print; it rests on only nine directional
decisions, so a single reversed call would move the terminal value materially; and the threshold
was chosen on a thirty-month validation window, which is itself a source of selection bias. What
the exhibit does legitimately show is that the same posterior can produce very different, and
more useful, behaviour when it is read through an economic loss function rather than a symmetric
classification rule. In that sense, the choice of loss function deserves at least as much
attention as the choice of scoring metric for the network.

\subsection{Summary of the empirical findings}

\begin{enumerate}
\item The dissertation's original protocol does not transfer well to an independent sample. On
our data it produces 23.3 per cent accuracy, below chance, because it combines non-persistent
regimes with a rolled-forward nowcast.
\item Redefining regimes on the magnitude of the monthly move, together with a prior that favours
persistence and strictly causal decoding, produces economically interpretable states and lifts
raw accuracy to 75.9 per cent.
\item Even so, that accuracy is statistically indistinguishable from a simple persistence rule
($p = 1.00$), and the balanced accuracy of 41.7 per cent reveals a model that leans too heavily
on the ``stagnant'' answer.
\item Turning the posterior into a trading position via an expected-return rule raises balanced
accuracy to 45.2 per cent and yields directional calls that are correct two times out of three,
at the cost of some raw accuracy.
\item Over a sample of only thirty months, none of these differences is statistically
significant. The honest reading is that the belief network may be picking up a small amount of
information about the direction of crude oil, but the sample is too short to demonstrate
conclusively that the signal is real.
\end{enumerate}

%==========================================================================
\section{Assessment of the contributions claimed by the dissertation}

The dissertation closes with eight claimed contributions (Alvi 66--67). Each one is assessed here
against the evidence the document itself provides, and against what came out of reproducing the
work. Our verdicts are pulled together in Table~\ref{tab:contrib} and then argued in detail
below.

\begin{table}[H]
\centering
\caption{Assessment of the eight claimed contributions of Alvi (2018).}
\label{tab:contrib}
\renewcommand{\arraystretch}{1.25}
\begin{tabular}{@{}p{0.5cm}p{7.4cm}p{2.6cm}p{4.2cm}@{}}
\toprule
\textbf{\#} & \textbf{Claim} & \textbf{Verdict} & \textbf{Principal evidence} \\
\midrule
1 & Bayesian views replace EGARCH-M views in a Black-Litterman model & Not achieved &
pp. 66--67; no allocation implemented \\
2 & Time series discretised by hidden Markov models as inputs to belief networks & Achieved &
Secs. 3.3, 4.2--4.2.5, pp. 42--43, 54--58 \\
3 & A working trading mechanism capable of independent decisions & Partially achieved &
Sec. 4.3.1, pp. 64--65 \\
4 & An autonomous system requiring no prior expert knowledge & Partially achieved &
Sec. 4.2.6, p. 59 \\
5 & The abstraction of the Python libraries clarifies the design process & Achieved &
Secs. 2.2.1--2.2.2, pp. 28--29 \\
6 & A global macro strategy returning more than high-frequency or fixed-income funds & Not achieved
& Sec. 4.3.1, p. 65 \\
7 & Better models of energy markets to inform policy & Partially achieved &
Sec. 3.4, pp. 43--45; figure p. 60 \\
8 & Amalgamation of existing research to increase alpha for commodity traders & Partially achieved
& Ch. 5, p. 67 \\
\bottomrule
\end{tabular}
\end{table}

\paragraph{1. Replacing EGARCH-M derived views with Bayesian model derived views in the
Black-Litterman framework (p. 66).} The claim picks up Beach and Orlov, whose 2007 paper derives
Black-Litterman views from an EGARCH-M model. The dissertation proposes the substitution in its
concluding section and cites the source, but no Black-Litterman allocation is ever actually
computed anywhere in the document. There is no prior equilibrium return vector, no view matrix,
no confidence matrix and no resulting portfolio. Chapter 4 ends with a single-asset long/short
rule on one barrel of crude (pp. 64--65), which is not a portfolio allocation. This contribution
is judged not achieved; it is a research proposal rather than a delivered result. Had it been
carried through, it would probably have been the most valuable of the eight, since a probabilistic
model that outputs a full posterior is exactly the kind of object that can supply the view
confidences Black-Litterman requires, and which practitioners usually have to set by hand.

\paragraph{2. Using time series discretised by hidden Markov models as inputs to belief networks
(p. 66).} This is the methodological core of the work and it is genuinely implemented. Section 3.3
lays out the design (pp. 42--43), Sections 4.2 to 4.2.5 fit the models and decode the regimes
(pp. 54--58), the regime-switch plot of the spot price on p. 56 shows the output, and the
discretised frames then feed the structure search. We reproduced the mechanism successfully,
which is about the strongest evidence one can have that it actually works. It is judged achieved,
with one qualification on novelty. Regime-switching representations of commodity markets long
predate the dissertation, so the specific contribution is the coupling of the decoded states to a
structure-learned belief network rather than the discretisation itself. That distinction matters,
because the coupling is the one part of the paper a practitioner could adopt directly, and our
own results show that its value depends entirely on how the emissions are defined.

\paragraph{3. A working trading mechanism capable of independent decision making (p. 66).}
Section 4.3.1 implements a rule that goes long on a forecast bull regime, short on a forecast bear
regime and flat otherwise, and plots the resulting value of one barrel against the price and
against the EIA Short Term Energy Outlook (figure on p. 65; United States, Energy Information
Administration). So the mechanism exists, and it does make decisions without human intervention.
It is not, however, a trading system in any operational sense: transaction costs, financing,
slippage and position limits are all absent, the evaluation covers only twenty-eight months, and
no risk-adjusted measure is computed. Our own equivalent exhibit (Figure~\ref{fig:strategy}) has
the same limitations, which is why it is presented here as an illustration. Partially achieved.

\paragraph{4. An autonomous system requiring no prior expert knowledge beyond dataset selection
(p. 66).} The dissertation's own code contradicts the claim in its strong form. On p. 59 the hill
climbing is initialised from an expert graph built from EIA relationships: six edges linking
OPEC and non-OPEC production, and OECD and non-OECD demand, to the spot price, and the search
then attaches the remaining variables. A search seeded with an expert structure is plainly not
free of expert knowledge. The seed determines which local optimum the greedy search lands in. Our
experiments make the point concrete: seeded and unseeded searches returned different networks
under the same score, and on the improved specification it was the unseeded sparse network that
was selected on the validation window. The claim holds for the unseeded variant, which the
dissertation does not evaluate on its own. Partially achieved.

\paragraph{5. The level of abstraction of the Python modelling libraries clarifies the design
process (p. 66--67).} Sections 2.2.1 and 2.2.2 (pp. 28--29) do explain how the libraries express
the underlying theory, and the exposition is clear and pedagogically useful. The claim is accepted
as achieved, while noting that it is a presentational contribution about third-party software
rather than a scientific one, and that it has aged. The library used for the hidden Markov models
is no longer maintained, which forced both our Group Work Project 2 and this project to substitute
a different implementation. The episode is a useful reminder that reproducibility claims resting
on a specific software stack decay quite quickly.

\paragraph{6. A systematic, event-driven global macro strategy capable of generating higher returns
than funds based on high-frequency or fixed-income strategies (p. 67).} No comparison with any
fund, index or strategy of either type is made anywhere in the document. The only benchmark
shown is the EIA Short Term Energy Outlook (p. 65), which is a forecast rather than an investable
strategy, and the comparison covers only a handful of episodes in 2016 and 2017, discussed
qualitatively. The author's own closing sentence on that page concedes that ``further research is
necessary to statistically conclude with confidence if the model performs better than the EIA's
STEO''. This claim is judged not achieved, and is regarded as the weakest of the eight. It
compares asset classes and strategy families on the back of a twenty-eight-month, single-asset
backtest with no risk adjustment.

\paragraph{7. Better models of energy markets, allowing policy makers to understand the underlying
structure (p. 67).} Section 3.4 (pp. 43--45) and the learned graph on p. 60 do deliver a structure
that can be inspected and argued about, which is more than a black-box forecaster offers, and
that is the real appeal of the graphical approach for a policy audience. The qualification is
stability. A structure meant to inform policy has to be robust, and ours was not: across eighteen
configurations the number of edges ranged from two to twenty-five, and the parents of the forecast
node moved around with the scoring function. Greedy search over a super-exponential space on 139
observations was never likely to recover a stable causal graph, and the dissertation does not
report any stability analysis such as bootstrap edge frequencies. Partially achieved; the
framework is appropriate, the evidential standard is not.

\paragraph{8. Amalgamating research from several disciplines and applying it to the commodity
markets to increase alpha for quantitative traders (p. 67).} The synthesis is real, and it is the
document's principal strength. Probabilistic graphical models, regime detection, macroeconomic
data engineering and a trading application are pulled together into a single coherent pipeline,
and the dissertation is a useful map of that territory. The clause about increasing alpha is
unsupported: no risk-adjusted return, benchmark comparison or significance test is reported
anywhere in the document. Partially achieved.

\subsection{Why the achieved contributions matter}

Two of the eight claims stand up to scrutiny, and both are worth keeping. The coupling of hidden
Markov discretisation with structure-learned belief networks (claim 2) matters because it gives
an energy desk something that neither a GARCH model nor a neural network naturally provides: a
forecast that comes with an explicit, inspectable statement of which variables it depended on.
When the model says the coming month is likely to be bearish, the Markov blanket says why, and a
risk committee can argue with the reasoning rather than just the number. The expository
contribution (claim 5) matters for a narrower but still practical reason: it lowers the cost of
entry for practitioners who need to audit such a model rather than build it from scratch.

The four partially achieved claims share a single shortcoming, and it is not conceptual but
evidential. The trading mechanism, the autonomy of the structure search, the policy relevance of
the graph and the alpha claim could all have been settled by evidence that the dissertation could
have produced but did not: transaction costs, a seeded-versus-unseeded comparison, bootstrap edge
stability, and a risk-adjusted benchmark. Our own results press the same point from the other
direction. Some of that evidence was produced here (exact confidence intervals, McNemar tests
against persistence and majority benchmarks, a coverage-and-hit-rate decomposition), and it
showed that what looked like a respectable 75.9 per cent accuracy is statistically
indistinguishable from simply repeating last month's answer. The lesson is that in this problem
the evidential apparatus is not a decoration bolted on after the modelling; it is what separates a
genuine result from a statistical artefact.

%==========================================================================
\section{Discussion}

\textit{This section is written for a reader who takes investment decisions and is deliberately
free of technical vocabulary.}

\subsection{What the study did and what it found}

Using information anyone can obtain free of charge (the price of oil, the strength of the dollar,
inflation, the interest rate set by the Federal Reserve and a measure of industrial activity),
we built a system that describes the oil market each month as being in one of three conditions:
falling, flat or rising. The system then works out the chance that each condition will hold in
the month ahead.

Tested against the two and a half years the system had never seen, it named the prevailing
condition correctly in roughly three months out of four. That number sounds good on paper, but it is
mostly misleading. The market was flat in most of those months, and a rule that simply always
said ``flat'' would have scored almost as well. The more honest measure is how the system behaved
at the moments that actually mattered. When it took a clear directional view, which happened in
roughly one month in three, it was right two times out of three. Acting on those views over
the period would have been materially more profitable than simply holding a barrel of oil the whole
time.

The original study on which this work is based was also repeated, and its method does not transfer
cleanly to our data. The reason is worth stating plainly, since it is probably the single most
useful thing we learned: the original method defined the market's condition purely by whether
the price went up or down last month. Defined in that way, the condition changes almost every
month, and a description that changes every month cannot reasonably support a decision held for
a quarter. Once the condition was instead defined by the size of the move, not just its
direction, the descriptions became stable, recognisable and usable. The system on its own
identified the 2014 collapse, the 2020 crash, the 2021 recovery and the 2025 correction as
clear episodes.

\subsection{The recommended course of action}

We do not recommend trading on this system alone. We do recommend using it in three specific ways.

\textbf{Use it to size risk, not to pick trades.} Its most reliable output is not a prediction
but a warning: the small number of months in which the evidence points clearly up or down.
In those months an energy book should be positioned more decisively. In the remaining two-thirds of
months, the sensible course is to reduce activity rather than invent a view.

\textbf{Use it as a challenge to the house view.} Because the system states which indicators drove
its conclusion, a committee can interrogate it. When it disagrees with the desk's opinion,
the disagreement pinpoints exactly which piece of evidence the two readings weigh differently.
That conversation is valuable even when the system turns out to be wrong.

\textbf{Keep any position modest and reviewed monthly.} The evidence base is thirty months and nine
directional calls. The apparent profitability excludes trading costs. On that basis, positions should
be small, should carry a predefined loss limit, and should be revisited each month as new figures
are published.

\subsection{The factors that move the result}

Three factors dominate, and a decision maker should keep an eye on them in this order.

\textbf{The definition of the market's condition.} This is the single largest influence on the
answer, larger than any refinement of the model downstream of it. A definition based on the size
of price moves produces conditions that persist for about four months. A definition based on
direction alone produces noise. Anyone commissioning a similar system should look closely at this
choice first.

\textbf{The dollar and the cost of money.} Oil is priced in dollars, so a stronger dollar makes it
more expensive everywhere else and weighs on demand. Higher policy rates raise the cost of carrying
inventory and financing production. Together they tend to push the market towards the falling condition,
and they were among the indicators the system attached to its forecast.

\textbf{Industrial activity.} It stands as a proxy for demand. When it weakens, the flat
condition tends to give way to the falling one.

A fourth factor is the reason for caution: events with no statistical precedent. Wars, production
decisions by exporting countries and the storage crisis of 2020 all dominate the price for
months at a time and are not in the information the system reads. It will always be late to
these, and no amount of statistical refinement will change that. It should be used as one
considered opinion among several, and never as the only one.

%==========================================================================
\section{Conclusion and practical takeaways}

The mini-capstone set out to decide whether probabilistic graphical models can forecast the
price of crude oil better than conventional time-series methods. Our answer, on sixteen years of
monthly data and a thirty-month hold-out window, is a qualified no, but a no with a useful
remainder.

\begin{enumerate}
\item \textbf{The framework is sound, the published implementation is fragile.} The dissertation's
testing accuracy could not be replicated, and the reason is now clear: its regime definition
produces non-persistent labels, and its evaluation protocol turns an accurate nowcast into a
systematically wrong forecast. Anyone picking this literature up in practice should test the
persistence of their regime labels before doing anything else.
\item \textbf{Discretisation matters more than model choice.} Changing the emission from the sign
of the monthly move to its magnitude, together with a prior that favours persistence, lifted
accuracy from 23.3 to 75.9 per cent. No change to the scoring function, in-degree limit or prior
inside the belief network came anywhere near that effect.
\item \textbf{Accuracy is the wrong headline number.} An accuracy of 75.9 per cent that is
indistinguishable from persistence ($p = 1.00$), and that comes with a balanced accuracy of 41.7
per cent, describes a model whose default answer is inaction. Balanced accuracy, coverage and hit
rate should be reported alongside accuracy in any regime-forecasting study.
\item \textbf{The loss function is part of the model.} Reading the same posterior with an
expected-return rule rather than a maximum a posteriori rule raised balanced accuracy by three and
a half points, produced directional calls that were right two times out of three, and changed the
economic outcome materially. This is the cheapest improvement available anywhere in the whole
pipeline.
\item \textbf{Sample size is the binding constraint.} Monthly macroeconomic data yield roughly two
hundred observations per decade and a half. Thirty-month hold-out windows produce confidence
intervals thirty points wide. Any future work in this direction should either move to a higher
frequency, pool across related commodities, or adopt a rolling-origin evaluation that re-uses the
sample honestly instead of relying on a single hold-out split (Hyndman and Athanasopoulos).
\end{enumerate}

For a risk manager, the practical value of the exercise lies less in the forecast itself than in
its transparency. A belief network states which variables its conclusion rests on. When it is
wrong, the reason is inspectable and the model can be corrected rather than merely retrained. In
a market whose structure shifts every few years, that property may matter more than a few points
of headline accuracy.

%==========================================================================
\newpage
\section*{Works Cited}
\addcontentsline{toc}{section}{Works Cited}

\noindent
Alvi, Danish A. \textit{Application of Probabilistic Graphical Models in Forecasting Crude Oil
Price}. 2018. University College London, dissertation. \textit{arXiv},
\url{https://arxiv.org/abs/1804.10869}.

\vspace{0.6em}\noindent
Ang, Andrew, and Allan Timmermann. ``Regime Changes and Financial Markets.'' \textit{Annual Review
of Financial Economics}, vol. 4, no. 1, 2012, pp. 313--337.

\vspace{0.6em}\noindent
Ankan, Ankur, and Abinash Panda. ``pgmpy: Probabilistic Graphical Models Using Python.''
\textit{Proceedings of the 14th Python in Science Conference}, 2015, pp. 6--11.

\vspace{0.6em}\noindent
Baumeister, Christiane, and Lutz Kilian. ``Forecasting the Real Price of Oil in a Changing World: A
Forecast Combination Approach.'' \textit{Journal of Business \& Economic Statistics}, vol. 33,
no. 3, 2015, pp. 338--351.

\vspace{0.6em}\noindent
Beach, Steven L., and Alexei G. Orlov. ``An Application of the Black--Litterman Model with
EGARCH-M-Derived Views for International Portfolio Management.'' \textit{Financial Markets and
Portfolio Management}, vol. 21, no. 2, 2007, pp. 147--166.

\vspace{0.6em}\noindent
Chickering, David Maxwell. ``Optimal Structure Identification with Greedy Search.'' \textit{Journal
of Machine Learning Research}, vol. 3, 2002, pp. 507--554.

\vspace{0.6em}\noindent
Cooper, Gregory F., and Edward Herskovits. ``A Bayesian Method for the Induction of Probabilistic
Networks from Data.'' \textit{Machine Learning}, vol. 9, no. 4, 1992, pp. 309--347.

\vspace{0.6em}\noindent
Fox, Emily B., et al. ``A Sticky HDP-HMM with Application to Speaker Diarization.'' \textit{The
Annals of Applied Statistics}, vol. 5, no. 2A, 2011, pp. 1020--1056.

\vspace{0.6em}\noindent
Hamilton, James D. ``A New Approach to the Economic Analysis of Nonstationary Time Series and the
Business Cycle.'' \textit{Econometrica}, vol. 57, no. 2, 1989, pp. 357--384.

\vspace{0.6em}\noindent
Heckerman, David, Dan Geiger, and David M. Chickering. ``Learning Bayesian Networks: The
Combination of Knowledge and Statistical Data.'' \textit{Machine Learning}, vol. 20, no. 3, 1995,
pp. 197--243.

\vspace{0.6em}\noindent
Hyndman, Rob J., and George Athanasopoulos. \textit{Forecasting: Principles and Practice}. 3rd ed.,
OTexts, 2021.

\vspace{0.6em}\noindent
Kilian, Lutz. ``Not All Oil Price Shocks Are Alike: Disentangling Demand and Supply Shocks in the
Crude Oil Market.'' \textit{American Economic Review}, vol. 99, no. 3, 2009, pp. 1053--1069.

\vspace{0.6em}\noindent
Koller, Daphne, and Nir Friedman. \textit{Probabilistic Graphical Models: Principles and
Techniques}. MIT Press, 2009.

\vspace{0.6em}\noindent
McNemar, Quinn. ``Note on the Sampling Error of the Difference between Correlated Proportions or
Percentages.'' \textit{Psychometrika}, vol. 12, no. 2, 1947, pp. 153--157.

\vspace{0.6em}\noindent
Rabiner, Lawrence R. ``A Tutorial on Hidden Markov Models and Selected Applications in Speech
Recognition.'' \textit{Proceedings of the IEEE}, vol. 77, no. 2, 1989, pp. 257--286.

\vspace{0.6em}\noindent
United States, Energy Information Administration. \textit{Short-Term Energy Outlook}. U.S.
Department of Energy, 2026, \url{https://www.eia.gov/outlooks/steo/}.

\vspace{0.6em}\noindent
United States, Federal Reserve Bank of St. Louis. \textit{Federal Reserve Economic Data (FRED)}.
2026, \url{https://fred.stlouisfed.org/}.

\end{document}
